% Part: first-order-logic
% Chapter: natural-deduction
% Section: quantifier-rules

\documentclass[../../../include/open-logic-section]{subfiles}

\begin{document}

\olfileid{fol}{ntd}{qrl}

\olsection{પરિમાણક નિયમો}

\subsection{$\lforall$ માટેના નિયમો}

\begin{defish}
\AxiomC{$!A(a)$}
\RightLabel{\Intro{\lforall}}
\UnaryInfC{$\lforall[x][\Atom{!A}{x}]$}
\DisplayProof
\hfill
\AxiomC{$\lforall[x][\Atom{!A}{x}]$}
\RightLabel{\Elim{\lforall}}
\UnaryInfC{$!A(t)$}
\DisplayProof
\end{defish}

$\lforall$ માટેના નિયમોમાં $t$ બંધ પદ છે (એવું પદ જેમાં કોઈ ચલ આવતો
નથી), અને $a$ એવો !!a{constant} છે જે ફલિતવાક્ય
$\lforall[x][!A(x)]$માં અથવા આધારવાક્ય~$!A(a)$ પર સમાપ્ત થતી
!!{derivation}માં !!{undischarged} રહેતી કોઈ ધારણામાં આવતો નથી.
અનુમાન \Intro{\lforall}ના $a$ને આપણે \emph{આઇગનચલ} કહીએ છીએ.\footnote{
ઉપરના નિયમમાં $a$ અચળ હોવા છતાં આપણે ``આઇગનચલ'' શબ્દ વાપરીએ છીએ.
તેના ઐતિહાસિક કારણો છે.}

\subsection{$\lexists$ માટેના નિયમો}

\begin{defish}
\AxiomC{$\Atom{!A}{t}$}
\RightLabel{\Intro{\lexists}}
\UnaryInfC{$\lexists[x][\Atom{!A}{x}]$}
\DisplayProof
\hfill
\AxiomC{$\lexists[x][\Atom{!A}{x}]$}
\AxiomC{[$\Atom{!A}{a}$]$^n$}
\DeduceC{$!C$}
\DischargeRule{\Elim{\lexists}}{n}
\BinaryInfC{$!C$}
\DisplayProof
\end{defish}

ફરીથી, $t$ બંધ પદ છે, અને $a$ એવો !!a{constant} છે જે આધારવાક્ય
$\lexists[x][!A(x)]$માં, ફલિતવાક્ય~$!C$માં, અથવા બંને આધારવાક્યો પર
સમાપ્ત થતી !!{derivation}sમાં !!{undischarged} રહેતી કોઈ ધારણામાં
(ધારણાઓ $!A(a)$ સિવાય) આવતો નથી. અનુમાન \Elim{\lexists}ના $a$ને
આપણે \emph{આઇગનચલ} કહીએ છીએ.

\Intro{\lforall} અથવા \Elim{\lexists} અનુમાન તરફ લઈ જતી
!!{derivation}sમાં કોઈ આઇગનચલ આધારવાક્યોમાં કે કોઈ !!{undischarged}
ધારણામાં ન આવે એવી શરતને \emph{આઇગનચલ-શરત} કહે છે.

\begin{explain}
યાદ કરો કે $!A$ એવા !!a{formula} હોય જેમાં !!{variable}~$x$ મુક્ત હોય,
તો આપણે તેને~$!A(x)$ લખીને દર્શાવીએ છીએ. એ જ સંદર્ભમાં $!A(t)$ એ
~$\Subst{!A}{t}{x}$નું સંક્ષિપ્ત રૂપ છે. તેથી $\Intro\lexists$ નિયમને
આ રીતે પણ લખી શકીએ:
\begin{prooftree}
\AxiomC{$\Subst{!A}{t}{x}$}
\RightLabel{\Intro{\lexists}}
\UnaryInfC{$\lexists[x][!A]$}
\end{prooftree}
નોંધો કે $t$ કદાચ $!A$માં પહેલેથી આવતું હોય; દાખલા તરીકે,
$!A$ કદાચ~$\Atom{\Obj P}{t,x}$ હોય. તેથી~$\Atom{\Obj P}{t,t}$ પરથી
$\lexists[x][\Atom{\Obj P}{t,x}]$નું અનુમાન કરવું એ
$\Intro\lexists$નો યોગ્ય ઉપયોગ છે---આધારવાક્યમાં આવતાં $t$નાં એક કે
વધુ, પણ જરૂરી નથી કે બધાં જ, આવર્તનોને બદ્ધ !!{variable}~$x$થી
``બદલી'' શકાય છે. પરંતુ $\Intro\lforall$ અને~$\Elim\lexists$ની
આઇગનચલ-શરતો મુજબ !!{constant}~$a$ એ~$!A$માં ન આવવો જોઈએ. તેથી
$\Atom{\Obj P}{a,a}$ પરથી~$\Intro\lforall$ વાપરીને
$\lforall[x][\Atom{\Obj P}{a,x}]$નું યોગ્ય અનુમાન કરી શકાતું નથી.
\end{explain}

\begin{explain}
\Intro{\lexists} અને \Elim{\lforall}માં કોઈ નિયંત્રણ નથી અને
પદ~$t$ કંઈ પણ હોઈ શકે છે, એટલે કોઈ શરતની ચિંતા કરવાની નથી. બીજી બાજુ,
\Elim{\lexists} અને \Intro{\lforall} નિયમોમાં આઇગનચલ-શરત મુજબ
!!{constant}~$a$ ફલિતવાક્યમાં કે કોઈ !!{undischarged} ધારણામાં ક્યાંય
ન આવવો જોઈએ. તંત્ર યથાર્થ રહે---અર્થાત્ જે !!{sentence}s જે
!!{undischarged} ધારણાઓમાંથી ફલિત થાય માત્ર તેમને જ !!{derive}s કરે---તે
માટે આ શરત જરૂરી છે. આ શરત વિના નીચેનું માન્ય ગણાત:
\begin{prooftree}
  \AxiomC{$\lexists[x][!A(x)]$}
  \AxiomC{$\Discharge{!A(a)}{1}$}
  \RightLabel{*\Intro{\lforall}}
  \UnaryInfC{$\lforall[x][!A(x)]$}
  \RightLabel{\Elim{\lexists}}
  \BinaryInfC{$\lforall[x][!A(x)]$}
\end{prooftree}
પરંતુ $\lexists[x][!A(x)] \Entails/ \lforall[x][!A(x)]$.

પરિમાણકો માટેના નિવારણ-નિયમો !!{variable}sના સ્થાને માત્ર બંધ પદો
મૂકવાની છૂટ આપે છે; તેથી !!{sentence}sના ગણમાંથી નિષ્પન્ન કરી શકાય એવું
કોઈ પણ !!{formula} પોતે પણ !!a{sentence} હોય છે.
\end{explain}


\end{document}
